%%%%%%%%%%%%%%%%%%%%%%%%%%%%%%%%%%%%%%%%%%
% Short Sectioned Assignment LaTeX Template Version 1.0 (5/5/12)
% This template has been downloaded from: http://www.LaTeXTemplates.com
% Original author:  Frits Wenneker (http://www.howtotex.com)
% License: CC BY-NC-SA 3.0 (http://creativecommons.org/licenses/by-nc-sa/3.0/)
%%%%%%%%%%%%%%%%%%%%%%%%%%%%%%%%%%%%%%%%%

%----------------------------------------------------------------------------------------
%	PACKAGES AND OTHER DOCUMENT CONFIGURATIONS
%----------------------------------------------------------------------------------------

\documentclass[paper=a4, fontsize=11pt]{scrartcl} % A4 paper and 11pt font size

% ---- Entrada y salida de texto -----

\usepackage[T1]{fontenc} % Use 8-bit encoding that has 256 glyphs
\usepackage[utf8]{inputenc}
%\usepackage{fourier} % Use the Adobe Utopia font for the document - comment this line to return to the LaTeX default

% ---- Idioma --------

\usepackage[spanish, es-tabla]{babel} % Selecciona el español para palabras introducidas automáticamente, p.ej. "septiembre" en la fecha y especifica que se use la palabra Tabla en vez de Cuadro

% ---- Otros paquetes ----

\usepackage{url} % ,href} %para incluir URLs e hipervínculos dentro del texto (aunque hay que instalar href)
\usepackage{amsmath,amsfonts,amsthm} % Math packages
%\usepackage{graphics,graphicx, floatrow} %para incluir imágenes y notas en las imágenes
\usepackage{graphics,graphicx, float} %para incluir imágenes y colocarlas

% Para hacer tablas comlejas
%\usepackage{multirow}
%\usepackage{threeparttable}

%\usepackage{sectsty} % Allows customizing section commands
%\allsectionsfont{\centering \normalfont\scshape} % Make all sections centered, the default font and small caps

\usepackage{fancyhdr} % Custom headers and footers
\pagestyle{fancyplain} % Makes all pages in the document conform to the custom headers and footers
\fancyhead{} % No page header - if you want one, create it in the same way as the footers below
\fancyfoot[L]{} % Empty left footer
\fancyfoot[C]{} % Empty center footer
\fancyfoot[R]{\thepage} % Page numbering for right footer
\renewcommand{\headrulewidth}{0pt} % Remove header underlines
\renewcommand{\footrulewidth}{0pt} % Remove footer underlines
\setlength{\headheight}{13.6pt} % Customize the height of the header

\numberwithin{equation}{section} % Number equations within sections (i.e. 1.1, 1.2, 2.1, 2.2 instead of 1, 2, 3, 4)
\numberwithin{figure}{section} % Number figures within sections (i.e. 1.1, 1.2, 2.1, 2.2 instead of 1, 2, 3, 4)
\numberwithin{table}{section} % Number tables within sections (i.e. 1.1, 1.2, 2.1, 2.2 instead of 1, 2, 3, 4)

\setlength\parindent{0pt} % Removes all indentation from paragraphs - comment this line for an assignment with lots of text

\newcommand{\horrule}[1]{\rule{\linewidth}{#1}} % Create horizontal rule command with 1 argument of height






%----------------------------------------------------------------------------------------
%	TÍTULO Y DATOS DEL ALUMNO
%----------------------------------------------------------------------------------------

\title{	
\normalfont \normalsize 
\textsc{\textbf{Grundlagen der Wissensverarbeitung} \\ Computer Science \\ Universität Hamburg} \\ [25pt] % Your university, school and/or department name(s)
~\\
~\\
~\\
\horrule{0.5pt} \\[0.4cm] % Thin top horizontal rule
\Huge Übung 1: Search Space Design \\ % The assignment title
\horrule{2pt} \\[0.5cm] % Thick bottom horizontal rule
~\\
~\\
}

\author{Rafael Ruz Gómez\\Miguel Robles Urquiza} % Nombre y apellidos

\date{\normalsize 5 November 2017} % Incluye la fecha actual

%----------------------------------------------------------------------------------------
% DOCUMENTO
%----------------------------------------------------------------------------------------

\begin{document}

\maketitle % Muestra el Título

\begin{figure}
	\centering
	\includegraphics[scale=0.8]{logo_uni_hamburg.png}
\end{figure}

\newpage %inserta un salto de página




%----------------------------------------------------------------------------------------
%	Question 1
%----------------------------------------------------------------------------------------

\huge{ \textbf{Exercise 1.1}}\\


\large{ \textbf{A core capability a student of computer science should obtain is to characterise 
and to model an environment for applying an algorithm or in our case a knowledge-based method.
For each distinction give a reason why the difference might be important when designing AI applications for the given environment.
Environments can be characterized by the following concepts:}

	\begin{itemize}
		\item \large{ \textbf{Fully observable <=> partially observable}}
	 	\item \large{ \textbf{Discrete <=> continuous}}
	 	\item \large{ \textbf{Deterministic <=> stochastic}}
	\end{itemize}

\large{ \textbf{Try to identify problems that exists in environments of one type but not in the other. You can use simple examples of environments and tasks to explain these problems.}}
\newline

\begin{itemize}
\item Fully observable: Too much space reserved in memory 
\item Partially observable: The agent has to recompute the path each step it does and discovers new environment.

\item Discrete: No problem found  
\item Continuous: Does not have a finite number of states.


\item Deterministic: Everything is structured, we can not take random actions \\      
\item Stochastic: There is no chance to predict the nex state 
\end{itemize}

%----------------------------------------------------------------------------------------
%	Question 2
%----------------------------------------------------------------------------------------

\huge{ \textbf{Exercise 1.2}}\\

\begin{enumerate}
\item 
\large{ \textbf{Imagine you have to design a route planning application for public transport (German: Öffentlicher Nahverkehr). What would be the state-space in this application? And what would be the semantic of nodes and edges in the graphs representing the state-space for this application?}}\\

States space is form by differents bus stops and where the buses are(Environment state) and in which bus or stop is the person(Agent state)
The nodes of this space are the bus stops and the arcs are the paths between two bus stops.\\

\item
\large{ \textbf{You are given two jugs, a 4-liter jug and a 3-liter jug. Neither has any measuring markers on it. There is a pump that can be used to fill the jugs with water. How can you get exactly 2 liters of water into the 4-liter jug? Solve this riddle using your knowledge about searching gained so far.}}\\

\large{ \textbf{a) Develop a formal model of the problem and select appropriate methods to solveit.}}
\begin{itemize}
\item States\\
3-liter jug fill\\
3-liter jug empty\\
3-liter jug with 1 liter\\
3-liter jug with 2 liters\\
4-liter jug fill\\
4-liter jug empty\\
4-liter jug with 1 liter\\
4-liter jug with 2 liters\\
4-liter jug with 3 liters\\


\item 
Start: Both jugs are empty\\
Goal: 2 liter in 4-liter jug\\

\item Transitions between states\\
Fill 4-liter jug\\
Fill 3-liter jug\\
Empty 4-liter jug\\
Empty 3-liter jug\\
Fill 4-liter jug with 3-liter jug\\
Fill 3-liter jug with 4-liter jug \\
\newline
\end{itemize}

\large{ \textbf{b)Imagine instead of water you would have to measure 2 liters of expensive wine, so wasting any liquid is not an option. Is the riddle still solvable? Why/Why not?}}

No, because there is no solution if you delete the empty actions from the transitions between states.
\newpage
\end{enumerate}


%------------------------------------------------

%\bibliography{citas} %archivo citas.bib que contiene las entradas 
%\bibliographystyle{plain} % hay varias formas de citar

\end{document}

